\scanpage{9}
\mspage{4}
which we shall confine ourselves to in what follows -- we thus have canonical homomorphisms
\[
(6)\qquad B_K\unsure{-|\longrightarrow}B_{\mathbb Q}\qquad\Pi_K\longrightarrow\Pi_{\mathbb Q}
\]
which are made explicit, when one has chosen an object of $\Pi_K$ i.e. a $\bar K/K$, whence a $\overline{\mathbb Q}/\mathbb Q$, by a homomorphism of profinite groups
\[
(7)\qquad E_{\bar K/K}\longrightarrow\Gamma_{\overline{\mathbb Q}/\mathbb Q}.
\]
In what follows, we shall always regard $B_K$, $\Pi_K$ or $E_{\bar K/K}$ as equipped with this additional structure -- it is the ``\uline{arithmetic}'' \deleted{f}morphisms (of topoi, of groupoids, or of profinite groups), will dominate the situation. !

Particular interest attaches to the kernel of (7), which I denote by $\pi_{\bar K/K}$ -- ``\unsure{one|of it} calls it the ``geometric part'' of the Galois group $E_{\bar K/K}$, as opposed to the quotient
\[
E_{\bar K/K}/\pi_{\bar K/K}=\Gamma^{\unsure{\deleted{a}|a}}_{\bar K/K}\hookrightarrow\Gamma_{\overline{\mathbb Q}/\mathbb Q},
\]
\margin{We shall denote by $\Gamma=\Gamma^{\unsure{\deleted{a}|\deleted{can}|can}}_{\bar K/K}$ this ``arithmetic part''}
which I call its
\margin{8}
